The transition from traditional cellular architectures to cell-free massive multiple-input multiple-output (MIMO) represents a paradigm shift in the design of sixth-generation (6G) wireless networks. By deploying a large number of distributed access points (APs) that cooperatively serve a smaller set of user equipments (UEs) within the same time-frequency resource, cell-free MIMO effectively eliminates cell-edge boundaries and provides substantial spatial diversity and macro-diversity gains. Unlike conventional co-located massive MIMO, the distributed nature of cell-free systems ensures that UEs are consistently in close proximity to multiple APs, thereby enhancing link reliability and significantly increasing the area spectral efficiency.

Despite these theoretical advantages, the practical implementation of cell-free MIMO is constrained by the "fronthaul bottleneck." In a typical deployment, APs must transmit either raw IQ samples or local signal estimates to a central processing unit (CPU) for joint detection and interference cancellation. The bandwidth of these fronthaul and midhaul links is often limited and subject to dynamic fluctuations due to network congestion or hardware constraints. Traditional quantization methods, such as Lloyd-Max quantization or fixed-rate scalar quantization, are generally designed to minimize mean squared error (MSE) based on static signal statistics. However, these approaches fail to adapt to real-time channel state information (CSI) and the specific requirements of the downstream detection task. Consequently, fixed quantization often leads to either inefficient utilization of available spectral resources in high-SNR regimes or severe degradation of detection accuracy in interference-limited scenarios.

Recent literature has explored deep learning (DL) to address the complexities of signal detection in massive MIMO. Model-driven DL architectures, such as those proposed in \cite{he2020model} and \cite{wei2020learned}, have demonstrated that unfolding iterative algorithms into neural networks can achieve near-optimal performance with reduced computational overhead. Furthermore, search-based methods like LISA \cite{sun2020learning} have shown robustness to imperfect CSI by learning optimal decision policies. In the context of distributed systems, deep joint source-channel coding (DJSCC) has emerged as a promising paradigm for task-oriented compression. However, a technical gap remains: existing graph neural network (GNN)-based detectors and DJSCC frameworks typically lack a mechanism for adaptive, differentiable bit-allocation that can prioritize information from critical APs while satisfying a global backhaul constraint. Most current distributed detection schemes assume symmetric quantization across all APs, which is suboptimal under heterogeneous channel conditions.

To address these challenges, this paper introduces IB-CellFree-GNN, a task-oriented successive refinement framework that integrates variational information bottleneck (VIB) principles with GNNs. The proposed framework enables APs to adaptively partition observations into progressive bit layers, allowing the CPU to perform resource-inference based on instantaneous network states. The primary contributions of this work are as follows:
\begin{itemize}
    \item We develop a successive refinement quantizer based on Learned Step-size Quantization (LSQ), which decomposes local linear minimum mean squared error (LMMSE) estimates at each AP into hierarchical bit-streams. This facilitates flexible bit-slicing where critical information is prioritized under stringent bandwidth limits.
    \item We propose a Gumbel-Softmax-based policy network for resource-inference. By transforming discrete bit-width selection into a differentiable operation during training, the network learns to autonomously allocate bits to APs with favorable channel conditions, maximizing task-relevant information transfer.
    \item We design a bit-aware GNN for CPU-side detection that utilizes a spatial attention mechanism. This architecture explicitly accounts for varying quantization noise levels across APs, allowing the GNN to compensate for low-precision signals from peripheral APs using high-precision features from core APs.
    \item Through numerical evaluations, we demonstrate that IB-CellFree-GNN achieves a 67.3\% reduction in bit error rate (BER) compared to fixed 2-bit quantization in heavy-load scenarios (12 users, 16 APs). Furthermore, the system demonstrates high resilience, maintaining a BER below 5\% even when 50\% of the backhaul links experience complete failure.
\end{itemize}

The remainder of this paper is organized as follows. Section II defines the cell-free MIMO system model and the backhaul-constrained detection problem. Section III details the proposed IB-CellFree-GNN architecture, including the successive refinement quantizer and the resource-inference module. Section IV provides an analysis of computational complexity. Section V presents numerical results and performance evaluations. Finally, Section VI concludes the paper.